
\section{Funding the project does not yet determine its accounts}
\label{sim:funding-measurement-control}
The price calculation leaves two different questions. Can the parties deliver
the promised cash in the required order? If they can, which amounts should
their financial statements recognise between issue and conversion? The first
question has a complete answer under the supplied synthetic cash balances.
The second still depends on the contract, reporting basis and measurement
evidence. Keeping the questions separate prevents a feasible transaction from
silently becoming an unsupported accounting history.

At the financing date the new holder pays $17000/441$ to the company and the
old owner contributes $9460/441$. Their receipts sum exactly to the project
price 60. The company then pays 60 to the seller. The transfer records identify
both parties, the date and the same amount in the two legal books. In this
example the company starts with no available cash for the purchase, so both
receipts must precede payment to the seller. There are six possible orders of
these three transfers. The two orders with the seller paid last succeed; the
other four fail at the first payment that exceeds available company cash.
Positive final net cash is not a substitute for an affordable cash prefix.

The four-party reference retains the company, old owner, new holder and seller
books separately. It checks two account charts and both owner cash 22 and cash
exactly equal to the required contribution. A refused later transfer preserves
the accepted prefix. Counterparty cost-basis entries are explicit synthetic
premises. They demonstrate reciprocal transfers and funding feasibility, not
the appropriate measurement of a real convertible instrument or the seller's
tax and profit consequences.

The limitation becomes visible if the 5\% pricing discount is reused as an
accounting yield. Suppose one incorrectly treats the face 80 claim as a plain
two-year zero-coupon liability and discounts its entire face at that rate.
The resulting liability and residual allocation would be
\begin{equation}
 L_0^{\rm forced}=\frac{80}{(21/20)^2}=\frac{32000}{441},
 \qquad P-L_0^{\rm forced}=-\frac{5000}{147}.
 \label{sim:forced-conversion-measurement}
\end{equation}
This negative residual is a diagnostic of that attempted mapping. The plain
unconditional cash claim has different payoffs from the contingent instrument
used to obtain $P$. Conversely, simply accreting $P$ at 5\% over both intervals
gives $85/2$, not face 80. Neither calculation supplies a replacement accounting
policy. The rate and cash flows appropriate to a liability component require
their own evidence.

For the conventional issuer compound-instrument case, the inspected
AASB~132 text measures the comparable liability first and assigns the residual
to equity. Its application guidance refers to instruments of comparable credit
status, substantially the same cash flows and terms, without the conversion
option\citep[paragraphs 31--32 and AG31]{aasb132}. Those conditions are not
established by the economic quote. The typed conversion demonstration instead
starts with a supplied carrying amount $C$. Its balanced transfer from liability
to equity remains conditional until classification, intervening measurement
and the conversion treatment are supported.

\section{Issuer and holder classify the same contract differently}
\label{sim:conversion-authority-scope}
The same contract is a claim issued by the company and an investment held by
the creditor. Those perspectives need not have matching classification
components. The source research therefore begins with the complete settlement
terms: who can affect the conversion signal, whether either party can elect
cash settlement, the currency and share terms, any redemption provisions, and
agreements involving other members of a group. The example supplies none of
these facts merely by stating that 120 shares will be issued.

Under the inspected AASB~132 provisions, fixed-for-fixed analysis is subject
to the definitions and exceptions for qualifying own-equity instruments.
A genuine contingent liability settlement beyond both parties' control raises
the separate conditions in paragraph 25. A derivative with settlement choices
has its own rule in paragraph 26. The non-derivative compound-instrument
assessment in paragraphs 28--32 cannot be inferred from a fixed share count
alone\citep[paragraphs 15--16, 22--26 and 28--32]{aasb132}. These passages identify
the questions to answer; they do not classify the underspecified fixture.

The holder's analysis is different. For an in-scope financial-asset host,
AASB~9 assesses the hybrid as a whole. Its fixed-number convertible-bond
example links the return to issuer equity and fails the solely-principal-and-
interest condition\citep[paragraphs 4.3.2 and B4.1.14, Instrument F]{aasb9scope2026}.
Issuer liability/equity components therefore cannot simply be copied into the
holder's asset book. Initial fair value, transaction costs, subsequent
measurement and relevant exceptions remain separate dependencies. Likewise,
AASB~132 AG32 expressly addresses conversion at maturity; it alone does not
justify the fixture's 2028 conversion before its 2029 terminal date.

The resulting holder interest also raises a reporting-perimeter question.
Section~\ref{sim:conversion-control-scope} follows the separate-statement,
consolidation and business-acquisition alternatives, including the possibility
of control before legal conversion. Its repaired dependency analysis preserves
uncertain control histories rather than choosing whichever journal balances.
The retained source compilations do not yet establish the complete accounting
authority for the fixture's 2027 issue and 2028 conversion. Classification,
measurement, control history and applicable edition therefore remain linked
questions with distinct evidence requirements.

\section{Evidence over an interval is stronger than two snapshots}
\label{sim:control-interval-evidence}
To support uninterrupted control, an analyst needs evidence covering the
intervening period. An empty event list supplies no such evidence. The selected
interval reference requires an explicit assertion that a control status holds
throughout a specified half-open interval, together with its subject, reporting
basis, source identity and availability to each actor. This is a representation
of an asserted coverage premise. It does not certify that the assertion is true
or that the source collection is complete.

For example, let one source assert no control from 30 December to 31 December,
and another assert control from 31 December to 1 January. A half-open interval
includes its starting instant and excludes its ending instant, so the shared
boundary is assigned once. If both assertions are eligible for the selected
actor view, they conditionally support a gain at that boundary. If the second
source is missing, the second interval remains unresolved. If it conflicts
with another eligible assertion, the conflict remains visible rather than being
resolved by whichever record was loaded last.

Effective time and information time do different jobs. A correction can describe
control during December yet become available to an analyst in January and to
a manager only in February. The analyst's January view can use it; the manager's
January view cannot. A later available revision replaces the corresponding
earlier assertion only within its declared subject and interval. Unknown actor
availability remains unresolved. The selected profile refuses partial-interval
replacement, whose meaning would require an additional rule.

The implementation reuses the exact fixed-offset timestamp parser and partitions
the query at every eligible interval boundary. It retains every covered,
uncovered, unknown or conflicting piece. Equivalent timestamp offsets denote
the same instant, and adding a future correction leaves an earlier view's
control conclusions unchanged. The checked cases include all two-cell combinations
of control, no control and unknown control, as well as revision conflicts,
coverage gaps, missing ancestry and corrupted source identities.

This gives the simulator a precise way to carry unresolved control history
alongside otherwise balanced company books. It does not yet connect an admitted
real history to acquisition measurement, changing group membership or complete
financial statements. The next financial composition must supply those methods
and their dated authority while preserving the source and actor-view distinctions.
Neither a successful funding transfer nor a complete conditional interval
partition can supply the missing accounting evidence.
